\documentclass[]{interact}

\usepackage{amsmath,amssymb,amsthm}
\usepackage{graphicx}
\usepackage{booktabs}
\usepackage[numbers,sort&compress]{natbib}
\bibliographystyle{tfnlm}

\theoremstyle{plain}
\newtheorem{theorem}{Theorem}
\newtheorem{proposition}{Proposition}
\newtheorem{lemma}{Lemma}
\newtheorem{assumption}{Assumption}
\newtheorem{definition}{Definition}
\theoremstyle{remark}
\newtheorem{remark}{Remark}

\newcommand{\R}{\mathbb{R}}
\newcommand{\E}{\mathbb{E}}
\newcommand{\Prob}{\mathbb{P}}
\newcommand{\law}{\mathrm{law}}
\newcommand{\Ocal}{\mathcal{O}}
\newcommand{\Sp}{\mathbb{S}}
\newcommand{\weakto}{\rightsquigarrow}
\newcommand{\He}{\mathrm{He}}
\newcommand{\push}[1]{{#1}_{\#}}
\newcommand{\Gp}{\mathbb{W}}
\DeclareMathOperator*{\argmin}{arg\,min}

\usepackage{url}
\begin{document}

\articletype{ARTICLE}

\title{Two-sample testing up to a rotation: calibration, a projected
non-Gaussianity test, and a Hermite view of contrast choice}

\author{\name{J. Touboul}\affil{Independent researcher}}

\thanks{CONTACT J. Touboul. Email: jack\_touboul@hotmail.com}

\maketitle

\begin{abstract}
We test whether two whitened planar distributions are equal up to an unknown
rotation---the alignment question behind two-dimensional independent component
analysis. Our first message is a caution: aligning a two-sample statistic by a
rotation search and calibrating by permutation is anticonservative, rejecting a
genuine rotation far too often. We give a projected non-Gaussianity statistic with
its limit law and a consistent multiplier bootstrap, and prove naive resampling
miscalibrates. Among valid competitors an ICA-canonicalized energy test is
competitive or more powerful, so the projected test's value is guaranteed
calibration, not power. On real breast-cancer data the projected test is the only
competitor that stays level-valid, while naive energy and even energy-after-ICA
over-reject. A Hermite view explains contrast choice: whitening annihilates the
first three Hermite orders, so a small corrected bank is used. All results are for
$d=2$.
\end{abstract}

\begin{keywords}
rotation alignment; independent component analysis; projection pursuit; two-sample
testing; multiplier bootstrap; Hermite polynomials
\end{keywords}

\section{Introduction}

Many questions ask not whether two distributions are equal, but whether they are
equal up to a transformation one does not know. Two-dimensional independent
component analysis (ICA) asks whether two whitened planar laws agree up to a
rotation. This paper is about testing that alignment, and about a pitfall in doing
so.

\paragraph{A calibration pitfall.} The obvious recipe is to take a good two-sample
statistic, minimise it over rotations of one sample, and calibrate by permutation.
This is invalid. The rotation search biases the observed statistic downward in a
way the permutation replicates only partially, and the resulting test rejects a
true rotation far too often (Section~\ref{sec:valid}; over $40\%$ in our runs).
Testing equality up to a group action is not a cosmetic modification of testing
equality: the nuisance must be handled in the calibration.

\paragraph{When the question is nontrivial.} Whitening removes second-order
structure. If the two laws have anisotropic covariance, whitening aligns them to
their covariance eigenaxes and the rotation is already resolved; a plain two-sample
test then suffices. The alignment question is substantive only in the isotropic
regime, where covariance does not pin the rotation and only higher-order structure
can. We work there throughout.

\paragraph{What is actually tested.} Our statistic compares, after the best
alignment, the direction \emph{profiles} of a scalar contrast. A profile is a
projection summary and does not determine the law, so profile alignment is
\emph{implied by} but does not \emph{imply} the alignment null. The test is thus a
calibrated screening test for a profile-based surrogate; a small bank of contrasts
widens the surrogate but does not make it exact. We state this plainly rather than
claim a characterisation the statistic cannot deliver.

\paragraph{Contributions.} (i) A limit law for the aligned statistic with a
consistent multiplier-bootstrap calibration, and a proof that naive resampling is
anticonservative (Sections~\ref{sec:prob}--\ref{sec:stat}); (ii) a comparison that
sorts competitors into valid and invalid, honest that among the valid tests an
ICA-canonicalized energy test is competitive or better (Section~\ref{sec:valid});
(iii) a Hermite analysis of contrast choice (Section~\ref{sec:xi}). All results are
for $d=2$; Section~\ref{sec:disc} discusses general $d$.

\paragraph{Related work.} Two-sample equality is tested by the energy distance
\cite{szekely}, kernel MMD \cite{gretton}, and, for dependence, the copula tests of
R\'emillard and Scaillet \cite{remillard}; none targets equality up to a group
action. Projection indices of non-Gaussianity go back to projection pursuit
\cite{friedman,huber} and underlie ICA \cite{comon,hyvarinen}; identifiability of
the rotation is the ICA identifiability question \cite{comon}, with modern
estimation theory in \cite{chenbickel,samworth}. Our focus is a calibrated test of
the rotation's existence, not its estimation.

\section{Setup}
\label{sec:setup}

Let $X,Y$ be $\R^2$-valued with finite moments and identity covariance
(``whitened''; in the isotropic regime whitening is a standardisation that leaves
the rotation unresolved). We observe independent i.i.d.\ unpaired samples of size
$n$.

\begin{definition}[Alignment null]
With $\Ocal(2)$ the orthogonal group, $H_0:\ \law(Y)=\push{R}\law(X)$ for some
$R\in\Ocal(2)$ ($\push{R}$ = pushforward).
\end{definition}

Fix a contrast $\xi:\R\to\R$. For $a_\theta=(\cos\theta,\sin\theta)$ and a whitened
law $P$, the \emph{profile} is $k_P(\theta)=\E_{Z\sim P}[\xi(a_\theta^\top Z)]$. A
rotation acts on the index by a shift, a reflection by $\theta\mapsto-\theta$, so
under $H_0$ the profiles agree after a shift. Define the profile discrepancy
\[
D_\xi=\min_{\epsilon\in\{\pm1\}}\min_s\int_0^{2\pi}\big(k_Y(\theta)-
k_X(\epsilon\theta+s)\big)^2\,d\theta .
\]

\begin{remark}[Surrogate hypothesis]
\label{rem:surrogate}
$D_\xi$ tests $H_0^\xi:$ ``the $\xi$-profiles align'', not $H_0$ itself. Since
$H_0\Rightarrow H_0^\xi$ but not conversely, the test controls the level under
$H_0$ and is a screening test for it: it can miss alternatives whose profiles
happen to align (Section~\ref{sec:xi} makes this precise in Hermite terms). This is
the price of reducing a $2$-D comparison to a $1$-D profile; the ICA-canonicalized
energy test of Section~\ref{sec:valid}, which retains the full distribution, does
not pay it, at the cost of a canonicalisation heuristic.
\end{remark}

With empirical profiles $\hat k_X,\hat k_Y$ on a fixed grid
$\theta_1,\dots,\theta_M$, the statistic is
\begin{equation}
\label{eq:Tn}
T_n=n\cdot\min_{\epsilon\in\{\pm1\}}\min_{s}\ \frac1M\sum_{g=1}^M
\big(\hat k_Y(\theta_g)-\hat k_X(\epsilon\theta_g+s)\big)^2 .
\end{equation}

\section{Limit law}
\label{sec:prob}

\begin{assumption}[$\xi$-integrability and Donsker class]
\label{ass:xi}
$\xi$ measurable, $\sup_a\E[\xi(a^\top Z)^2]<\infty$ for both laws, and
$\{z\mapsto\xi(a^\top z):a\in\Sp^1\}$ Donsker for both.
\end{assumption}

\begin{lemma}[Functional CLT]
\label{lem:clt}
Under Assumption~\ref{ass:xi}, $\sqrt n(\hat k_X-k_X)\weakto\Gp_X$ and
$\sqrt n(\hat k_Y-k_Y)\weakto\Gp_Y$ independently, mean-zero Gaussian with
covariance $\Gamma_\xi^\bullet(\theta,\theta')=\mathrm{Cov}(\xi(a_\theta^\top Z),
\xi(a_{\theta'}^\top Z))$.
\end{lemma}

Let $s_0$ be the population-optimal shift ($D(s_0)=0$ under $H_0$).

\begin{assumption}[Non-degenerate alignment]
\label{ass:nondeg}
$s_0$ is isolated modulo reflection, $D''(s_0)>0$, $k_X'(\cdot+s_0)\not\equiv0$.
\end{assumption}

\begin{theorem}[Limit law]
\label{thm:limit}
Under $H_0$ and Assumptions~\ref{ass:xi}--\ref{ass:nondeg}, for a fixed grid of $M$
directions,
\[
T_n\weakto L_\xi^{(M)}:=\|\Delta\|^2-\frac{\langle\Delta,k_X'(\cdot+s_0)\rangle^2}
{\|k_X'(\cdot+s_0)\|^2},\quad \Delta=\Gp_Y-\Gp_X(\cdot+s_0),
\]
a finite weighted sum of independent $\chi^2_1$ variables (weights = eigenvalues of
the covariance of $\Delta$ off the score direction). As $M\to\infty$ the discrete
inner products converge to their integral analogues.
\end{theorem}

\begin{remark}[Fixed grid]
\label{rem:grid}
We treat $M$ as fixed; $T_n$ converges to the $M$-dimensional quadratic form
$L_\xi^{(M)}$, not to a continuous functional. This is deliberate: the bootstrap of
Section~\ref{sec:stat} uses the \emph{same} grid, so it reproduces $L_\xi^{(M)}$
for that $M$ and validity does not require $M\to\infty$. Larger $M$ refines the
alignment resolution at $O(M)$ cost.
\end{remark}

\begin{remark}[Role of the limit law]
$L_\xi^{(M)}$ is not a plug-in critical value---it depends on unknown moments. It
is the law the bootstrap reproduces, and under contiguous alternatives
$k_Y=k_X(\cdot+s_0)+h/\sqrt n$ it becomes the noncentral $\|\Delta+\Pi h\|^2$ ($\Pi$
= projection off the score), giving the local power.
\end{remark}

\section{Calibration and consistency}
\label{sec:stat}

\paragraph{Why naive calibration fails.} A permutation two-sample bootstrap tests
equality of laws, not equality up to rotation, and rejects true alignments; a
plug-in dropping the score term overstates the statistic; a split estimating the
alignment separately mis-centres the correction. Most consequentially
(Section~\ref{sec:valid}), \emph{minimising a two-sample statistic over rotations
and permuting} is anticonservative: the rotation search lowers the observed value
more than it lowers the permutation replicates.

\paragraph{Multiplier bootstrap.} Let $\hat\psi^X_g,\hat\psi^Y_g$ be the whitened
profile influence functions and $w_1,\dots,w_n$ (resp.\ $w_1',\dots$) i.i.d.\
mean-zero variance-one multipliers independent of the data. With
$\Gp_X^{*}(\theta_g)=n^{-1/2}\sum_i w_i\hat\psi^X_g(X_i)$ and $\Gp_Y^{*}$
analogously, form $T_n^{*}$ from $\Delta^{*}=\Gp_Y^{*}-\Gp_X^{*}(\cdot+\hat s)$
exactly as $T_n$ (the same aligned minimisation reproduces the score correction).

\begin{proposition}[Bootstrap consistency]
\label{prop:boot}
Under $H_0$ and Assumptions~\ref{ass:xi}--\ref{ass:nondeg}, conditionally on the
data and in probability $T_n^{*}\weakto L_\xi^{(M)}$; the test rejecting when $T_n$
exceeds the conditional $(1-\alpha)$ quantile has asymptotic level $\alpha$.
\end{proposition}

\begin{proposition}[Estimated whitening]
\label{prop:white}
With sample whitening, Lemma~\ref{lem:clt} holds with $\psi^X_g$ replaced by
$\bar\psi^X_g=\psi^X_g+\Lambda_g^\top\mathrm{vech}(XX^\top-\E XX^\top)$; the
bootstrap stays consistent if the multipliers weight $\bar\psi$.
\end{proposition}

\begin{proposition}[Consistency]
\label{prop:power}
If $D_\xi>0$ then $T_n\to\infty$ and the power tends to one. If $D_\xi=0$ while
$H_0$ fails (Remark~\ref{rem:surrogate}), the test is not consistent for that
contrast---the reason for the bank of Section~\ref{sec:xi}.
\end{proposition}

\paragraph{Numerical calibration.} Figure~\ref{fig:calib} shows the bootstrap in
the isotropic light-tailed regime: the level is controlled---indeed conservative at
these sizes---and power tends to one.

\begin{figure}[t]
\centering
\includegraphics[width=0.60\textwidth]{fig_calibration.png}
\caption{Calibration, isotropic light-tailed regime. Dashed: null level; solid:
power vs a Gaussian alternative; dotted: nominal $0.05$.}
\label{fig:calib}
\end{figure}

\section{Which tests are valid, and how they compare}
\label{sec:valid}

We compare five procedures at $n=200$ over $80$ replications
(Figure~\ref{fig:power}, Table~\ref{tab:power}): the projected test; an
ICA-canonicalized energy test (canonicalise each sample by FastICA, then an energy
two-sample test); a rotation-aligned energy test (minimise energy over rotations,
permutation-calibrated); a naive energy test; and a naive kernel (MMD) test.

\begin{figure}[t]
\centering
\includegraphics[width=0.72\textwidth]{fig_power_comparison.png}
\caption{Validity and power ($n=200$, $80$ replications, $\pm1$ s.e.). The two
naive aligned procedures over-reject the rotation null and are invalid; the
projected test and energy-after-ICA control the level. Among the valid tests,
energy-after-ICA is at least as powerful.}
\label{fig:power}
\end{figure}

\begin{table}[t]
\centering
\begin{tabular}{lccc}
\toprule
method & $H_0$ (rotation) & power (Gaussian) & power (uniform)\\
\midrule
projected test (ours)     & $0.00\pm.00$ & $0.81\pm.04$ & $0.04\pm.02$\\
energy after ICA          & $0.06\pm.03$ & $0.98\pm.02$ & $0.25\pm.05$\\
energy, rotation-aligned  & $0.54\pm.06$ & $1.00\pm.00$ & $1.00\pm.00$\\
energy, naive             & $0.84\pm.04$ & $0.99\pm.01$ & $0.53\pm.06$\\
MMD, naive                & $0.00\pm.00$ & $0.00\pm.00$ & $0.00\pm.00$\\
\bottomrule
\end{tabular}
\caption{Rejection rates, isotropic regime, $n=200$. The rotation-aligned and naive
energy tests are invalid; the ``power'' columns are meaningful only for the
level-controlling rows.}
\label{tab:power}
\end{table}

Two findings stand out. First, \emph{validity}: aligning a two-sample statistic by
a rotation search and permuting does not control the level---the rotation-aligned
energy test rejects the rotation null $54\%$ of the time, and the naive energy test
$84\%$. Only the projected test and energy-after-ICA are valid. Second,
\emph{power}: among the valid tests, energy-after-ICA is competitive and, on these
alternatives, more powerful than the projected test---unsurprising, since it retains
the full distribution after alignment while the projected test compresses to a
$1$-D profile. We therefore do not claim power superiority for the projected test.
Its case rests on Sections~\ref{sec:prob}--\ref{sec:stat}: an explicit limit law
and a bootstrap with a proven guarantee, where energy-after-ICA relies on a
canonicalisation heuristic tied to ICA identifiability.

\section{A Hermite view of contrast choice}
\label{sec:xi}

Expand $\xi$ in probabilist Hermite polynomials, $\xi=\sum_j c_j\He_j$,
$c_j=\E_{N(0,1)}[\xi\He_j]/j!$. For whitened $Z$, $\E[\He_1(a^\top Z)]=0$ (centred)
and $\E[\He_2(a^\top Z)]=\mathrm{Var}(a^\top Z)-1=0$ (identity covariance), so
\begin{equation}
\label{eq:hermite}
k_P(\theta)=c_0(\xi)+\sum_{j\ge3}c_j(\xi)\,m_j^P(\theta),\qquad
m_j^P(\theta)=\E[\He_j(a_\theta^\top Z)] :
\end{equation}
\emph{whitening annihilates Hermite orders $0,1,2$}, and the test sees only the
supra-quadratic part of $\xi$. Three consequences.

\emph{Non-degeneracy} (Assumption~\ref{ass:nondeg}) holds iff $\sum_{j\ge3}c_j
m_j^X$ is non-constant---the null must carry directional structure at some order
$j\ge3$ with $c_j\ne0$. This is the surrogate gap of Remark~\ref{rem:surrogate} in
coordinates, and links to ICA identifiability \cite{comon}.

\emph{No single contrast dominates.} A contrast is blind to any order with
$c_j=0$: $u^4$ (mass at $j{=}4$) is blind to skewness, an odd contrast to symmetric
kurtosis. Figure~\ref{fig:diss} and Table~\ref{tab:diss} show the dissociation:
$u^3$ detects a skewed alternative but not a Gaussian one, $u^4$ the reverse. All
contrasts with supra-quadratic support are consistent (Table~\ref{tab:hermite}),
differing only in efficiency; $\log\cosh$, standard in ICA, is the least efficient,
its Hermite mass sitting at the annihilated order two.

\emph{Combine with a corrected bank.} Since no fixed contrast is uniformly best,
combine a few. Pooling the discrepancies into one statistic is delicate---the sum
is dominated by the noisiest order and depends on scaling---so we use a scale-free
multiplicity-corrected bank: reject if $\min_j p_j<\alpha/|\text{bank}|$. With
$\{u^3,u^4\}$ this dominates either contrast alone (Table~\ref{tab:diss}); the
Bonferroni correction is conservative because the $p_j$ are positively correlated,
and a jointly calibrated combination would recover a little power.

\begin{figure}[t]
\centering
\includegraphics[width=0.62\textwidth]{fig_dissociation.png}
\caption{No single contrast dominates. $u^4$ misses the skewness alternative, $u^3$
the kurtosis one; the corrected bank $\{u^3,u^4\}$ handles both.}
\label{fig:diss}
\end{figure}

\begin{table}[t]
\centering
\begin{tabular}{lcccc}
\toprule
contrast & $c_4$ & $c_6$ & supra-quad.\ variance & data factor vs $u^4$\\
\midrule
$u^4$     & $1.00$   & $0$      & $24.0$   & $1\times$\\
$u^6$     & $15.0$   & $1.0$    & $6120$   & $0.004\times$\\
$|u|$     & $-0.033$ & $0.003$  & $0.042$  & $574\times$\\
$\log\cosh$ & $-0.015$ & $0.001$ & $0.006$ & $3793\times$\\
\bottomrule
\end{tabular}
\caption{Post-whitening Hermite content. All contrasts are consistent; the data
factor is the relative sample size for comparable signal. Efficiency is not
monotone in $|c_j|$: the bounded $|u|$ has tiny coefficients but tiny noise.}
\label{tab:hermite}
\end{table}

\begin{table}[t]
\centering
\begin{tabular}{lccc}
\toprule
contrast/method & $H_0$ & power (kurtosis alt.) & power (skewness alt.)\\
\midrule
$u^3$ (skewness)          & 0.02 & 0.00 & 0.95\\
$u^4$ (kurtosis)          & 0.00 & 1.00 & 0.77\\
bank $\{u^3,u^4\}$ (Bonf.)& 0.00 & 0.97 & 0.92\\
\bottomrule
\end{tabular}
\caption{Contrast choice, $n=400$. Single contrasts dissociate; the bank handles
both.}
\label{tab:diss}
\end{table}

\section{A real-data illustration}
\label{sec:realdata}

The regime this paper targets---whitened, isotropic, non-Gaussian---occurs in
real bivariate data, and the calibration caution of Section~\ref{sec:valid} is not
an artifact of Gaussian simulation. We illustrate on the Wisconsin diagnostic
breast-cancer data, taking two skewed real features (mean area and mean concavity;
sample skewness $1.6$ and $1.4$).

\emph{Calibration on real marginals.} To obtain a rotation null that is true by
construction while keeping genuinely real marginals, we split the sample into two
halves and apply an independent random rotation to each; the two halves are then
equal up to a rotation. Over $100$ such splits, Table~\ref{tab:real} reports the
rejection rate at level $0.05$. The naive energy test rejects the true rotation in
$76\%$--$83\%$ of splits---the over-rejection of Section~\ref{sec:valid}, now on
real data. The projected test controls the level (no rejections). Strikingly, the
ICA-canonicalized energy test---valid, and often more powerful, in the simulations
of Section~\ref{sec:valid}---over-rejects here ($33\%$--$64\%$): on these marginals
the two whitened components carry similar higher-order structure, so the ICA
canonicalization resolves the two halves to different orientations and the
comparison ceases to be level-valid. The rotation-aligned energy test, by contrast,
controls the level on these data ($0.00$--$0.04$), so its over-rejection in the
isotropic simulation of Section~\ref{sec:valid} is regime-dependent; the naive energy
test, however, over-rejects in both settings, and the projected test controls the
level in both. The projected test's calibration does not depend on an ICA fit
converging, and this is where that guarantee earns its place.

\begin{table}[htbp]
\centering
\begin{tabular}{lcc}
\toprule
rejection under the rotation null ($\alpha=0.05$) & area / concavity & radius / worst-conc. \\
\midrule
projected test (ours)      & \textbf{0.00} & \textbf{0.00} \\
energy after ICA           & 0.33 & 0.64 \\
energy, rotation-aligned   & 0.00 & 0.04 \\
energy, naive              & 0.76 & 0.83 \\
\bottomrule
\end{tabular}
\caption{Calibration on real marginals (breast-cancer), two feature pairs. The
projected test controls the level; the naive energy test grossly over-rejects; and
the ICA-canonicalized energy test, valid in simulation, loses level control on real
data.}
\label{tab:real}
\end{table}

\begin{figure}[htbp]
\centering
\includegraphics[width=0.68\textwidth]{fig_realdata.png}
\caption{Rejection rate under a true rotation null on real marginals. Only the
projected test controls the level across both feature pairs; naive energy and
energy-after-ICA over-reject.}
\label{fig:real}
\end{figure}

\emph{A two-sample comparison.} Testing the benign against the malignant class on
the same two features, the projected test does not reject equality up to a rotation
($p=0.45$), while the naive and ICA tests return $p=0.17$ and $p=0.03$. Given the
calibration above---where the naive and ICA tests reject a \emph{true} null at high
rates---their $p$-values here cannot be read as evidence; only the level-valid
projected test supports an interpretable conclusion. The lesson mirrors the paper's
thesis: on real data the projected test is the one procedure that stays valid, so
its verdict is the one a practitioner can trust.

\section{Discussion}
\label{sec:disc}

The practical message is twofold: testing equality up to a rotation demands
calibration designed for the nuisance---naive aligned two-sample tests are
anticonservative---and, once validity is secured, an ICA-canonicalized energy test
is a strong baseline that the projected test does not beat on power. The projected
test's value is its explicit theory and guaranteed calibration, together with the
Hermite reading of contrast choice.

The honest limits are several. The theory is for $d=2$; general $d$ replaces the
scalar shift by an $\argmin$ over the $\Ocal(d)$ manifold with $d(d-1)/2$ nuisance
directions and can have non-isolated minima when components share higher moments,
so Theorem~\ref{thm:limit} does not extend without further work; the same
degeneracy makes Assumption~\ref{ass:nondeg} fail, and, tellingly, it is exactly
when the alignment is ambiguous that ICA canonicalisation also fails, so no method
gains there. The method targets a profile surrogate (Remark~\ref{rem:surrogate}),
mitigated but not removed by the bank. It suits the isotropic regime; under
anisotropic covariance whitening already resolves the rotation. The kurtosis
contrast needs eighth moments, so heavy-tailed sources require large samples. A
robust divergence analogue for side-channel analysis is a natural applied direction
requiring validation on real data, left to future work.

\section*{Disclosure statement}
The author reports no competing interests.

\section*{Funding}
No funding was received.

\section*{Data availability statement}
The real-data illustration uses the Wisconsin Diagnostic Breast Cancer dataset, a
public benchmark available in the \texttt{scikit-learn} library and originally from
the UCI Machine Learning Repository. No new data were generated. All results---the
simulation studies and the real-data analysis---are fully reproducible from the code
provided as supplementary material and archived at
\url{https://doi.org/10.5281/zenodo.23023835}, which runs with a fixed random seed.

\appendix

\section{Reminders and proofs}
\label{app}

\emph{Donsker/Hermite.} $\{\theta\mapsto\xi(a_\theta^\top z)\}$ is Donsker under a
square-integrable envelope ($\|z\|^4$ for $u^4$, needing $\E\|Z\|^8<\infty$;
$O(\|z\|)$ for $\log\cosh$). The probabilist polynomials satisfy
$\E_{N(0,1)}[\He_j\He_k]=j!\delta_{jk}$, giving \eqref{eq:hermite}; orders $0,1,2$
integrate to constants under centring and identity covariance.

\medskip\noindent\emph{Theorem~\ref{thm:limit}.} Let $D_n(s)=\frac1M\sum_g(\hat
k_Y(\theta_g)-\hat k_X(\theta_g+s))^2$. At $s_0$, $nD_n(s_0)\weakto\|\Delta\|^2$ by
Lemma~\ref{lem:clt}. By Assumption~\ref{ass:nondeg} and the argmax theorem
$\hat s-s_0=O_\Prob(n^{-1/2})$; a second-order expansion with $D_n'(s_0)=
O_\Prob(n^{-1/2})$, $D_n''(s_0)\to D''(s_0)$ gives $n\min_s D_n(s)=nD_n(s_0)-
\tfrac12\{\sqrt n D_n'(s_0)\}^2/D_n''(s_0)+o_\Prob(1)$. Since $\sqrt n D_n'(s_0)
\weakto-2\langle\Delta,k_X'(\cdot+s_0)\rangle$ and $D_n''(s_0)\to2\|k_X'(\cdot+s_0)
\|^2$, continuous mapping yields $L_\xi^{(M)}$; the reflection is a continuous
minimum over two branches.

\medskip\noindent\emph{Propositions~\ref{prop:boot}--\ref{prop:white}.} With
$W(\Sigma)$ smooth, a first-order expansion gives the augmented influence function
$\bar\psi^X_a=\psi^X_a+\Lambda_a^\top\mathrm{vech}(XX^\top-\Sigma)$. Conditionally,
$\Gp^{*}_\bullet$ sums independent mean-zero terms with conditional covariance
$\to\Gamma_\xi^\bullet$, so by the multiplier CLT $\Gp^{*}_\bullet\weakto
\Gp_\bullet$ conditionally in probability; the norm and score projection are
continuous and $\hat s\to s_0$, giving $T_n^{*}\weakto L_\xi^{(M)}$. The
anticonservatism of the rotation-aligned permutation test follows because the
observed statistic minimises over a continuum of rotations while each permutation
replicate minimises over the same continuum on exchangeable data.

\begin{thebibliography}{99}
\bibitem{szekely} G.~Sz\'ekely and M.~Rizzo. Energy statistics.
\emph{J.\ Statist.\ Plann.\ Inference}, 143:1249--1272, 2013.
\bibitem{gretton} A.~Gretton et al. A kernel two-sample test. \emph{JMLR},
13:723--773, 2012.
\bibitem{remillard} B.~R\'emillard and O.~Scaillet. Testing for equality between
two copulas. \emph{J.\ Multivariate Anal.}, 100:377--386, 2009.
\bibitem{friedman} J.~Friedman and J.~Tukey. A projection pursuit algorithm for
exploratory data analysis. \emph{IEEE Trans.\ Comput.}, C-23:881--890, 1974.
\bibitem{huber} P.~Huber. Projection pursuit. \emph{Ann.\ Statist.},
13:435--475, 1985.
\bibitem{comon} P.~Comon. Independent component analysis, a new concept?
\emph{Signal Processing}, 36:287--314, 1994.
\bibitem{hyvarinen} A.~Hyv\"arinen and E.~Oja. Independent component analysis:
algorithms and applications. \emph{Neural Networks}, 13:411--430, 2000.
\bibitem{chenbickel} A.~Chen and P.~Bickel. Efficient independent component
analysis. \emph{Ann.\ Statist.}, 34:2825--2855, 2006.
\bibitem{samworth} R.~Samworth and M.~Yuan. Independent component analysis via
nonparametric maximum likelihood estimation. \emph{Ann.\ Statist.},
40:2973--3002, 2012.
\bibitem{vdvw} A.~van der Vaart and J.~Wellner. \emph{Weak Convergence and
Empirical Processes}. Springer, 1996.
\end{thebibliography}

\end{document}
